\naslov{Nenadzorovano učenje}

Pri nenadzorovanem učenju nimamo ciljne spremenljivke $Y$.
Namesto napovedi iščemo povezane skupine podatkov (\enquote{clustering}).
Kot vhodni podatek dobimo množico podatkov $S$ z metriko $d$, in želeno število
skupin $k$.
Prva možnost tu je algoritem HAC.
Začnemo s $\abs{S}$ enojci, dokler imamo več kot eno skupino, najbližji dve
združimo.
Rezultat je dendrogram, kjer $y$ koordinata povezave med skupinama prikaže
medsebojno razdaljo.
Razdaljo med skupinama $C_1$ in $C_2$ lahko definiramo na več načinov, ali kot
najmanjšo/največjo/povprečno razdaljo med primeri, ali pa bolj ekonomično; za
tretjo skupino $C$ je
\[
  D(C_1 \unija C_2, C) = \alpha_1 D(C_1, C) + \alpha_2 D(C_2, C) + \beta D(C_1,
  C_2)
\]
za
\begin{gather*}
  \alpha_i = \frac{\abs{C_i} + \abs{C}}{\abs{C_1} + \abs{C_2} + \abs{C}}, \\
  \beta = - \frac{\abs{C}}{\abs{C_1} + \abs{C_2} + \abs{C}}.
\end{gather*}
Temu pravimo \pojem{Wardova formula}.

\vprasanje{Opiši delovanje algoritma HAC in povej Wardovo formulo.}

Druga metoda je $k$-means clustering.
Tukaj iščemo
\[
  \min_{\abs{\mathcal{C}} = k} \sum_{C \in \mathcal{C}} \abs{C} \var(C),
\]
kar lahko naredimo za evklidsko razdaljo.
Izračunamo \pojem{centroid} (težišče)
\[
  c = \operatorname{arg\,min}_u \sum_{v \in S} d(u,v)
  = \inv{\abs{S}} \sum_{v \in S} \vec{v},
\]
kar dobimo z računanjem odvodov.
Varianca skupine je vsota kvadratov razdalj med primeri in centroidom.

Razbitje izračunamo iterativno.
Začnemo z naključno razporeditvijo, in potem iterativno izračunamo centroide
vseh skupin ter vse primere povežemo s tistim centroidom, ki jim je najbližje.
Iteracijo ustavimo, ko ni več spremembe.

Če imamo neevklidsko mero razdalje, se lahko zgodi, da centroida ne znamo
poiskati.
V tem primeru ga nadomestimo z \pojem{medoidom}
\[
  m = \operatorname{arg\,min}_{u \in S} \sum_{v \in S} d(u,v).
\]

\vprasanje{Opiši delovanje algoritma $k$-means clustering. Kaj naredimo, če
  nimamo evklidske metrike?}

Kvaliteto razvrščanja merimo s povezanostjo in ločenostjo primerov od ostalih.
\pojem{Povezanost} je definirana kot
\[
  a(u) = \inv{\abs{C_u} - 1} \sum_{v \in C_u} d(u,v),
\]
kjer je $C_u$ skupina, kateri pripada $u$.
Delimo s $\abs{C_u} - 1$, ker je razdalja $d(u,u) = 0$.
\pojem{Ločenost} pa je definirana kot
\[
  b(u) = \min_{C \in \mathcal{C}, C \ne C_u} \inv{\abs{C}} \sum_{v \in C} d(u,v).
\]
Ti vrednosti kombiniramo v \pojem{obris}
\[
  s(u) = \frac{b(u) - a(u)}{\max(a(u), b(u))},
\]
ki ima vrednosti med $-1$ in $1$.
Primer je razvrščen dobro za $s$ blizu $1$ in slabo za $s$ blizu $-1$; mera
kvalitete celotnega razvrščanja bo povprečje $s(u)$ čez vse $u$.

\vprasanje{Kako vrednotimo razvrščanje?}

% LocalWords:  HAC dendrogram Wardova Wardovo means clustering centroidom
% LocalWords:  centroide centroida medoidom
